\documentclass[a4paper,landscape,12pt]{extarticle}
\usepackage[left=0.5in,right=0.5in,top=0.4in,bottom=0.55in]{geometry}

\usepackage{graphicx}
\usepackage[normalem]{ulem}
\usepackage{multicol}
\usepackage{array}
\usepackage{caption}
\usepackage{fancyhdr}
\usepackage{lastpage}
\usepackage{amsmath,amssymb}
\usepackage{booktabs}
\usepackage{listings}
\usepackage{xcolor}
\usepackage{hyperref}

% ---------- Listings (code blocks) ----------
\lstset{
  language=Python,
  basicstyle=\ttfamily\Large,
  keywordstyle=\color{blue}\bfseries,
  commentstyle=\color{gray},
  stringstyle=\color{teal},
  numbers=left,
  numberstyle=\Large\color{gray},
  frame=single,
  breaklines=true,
  breakatwhitespace=true,
  columns=fullflexible,
  keepspaces=true,
  showstringspaces=false
}

\author{Jacky Li}

% Page style with slide numbers
\pagestyle{fancy}
\fancyhf{}
\renewcommand{\headrulewidth}{0pt}
\fancyfoot[R]{\Large\thepage\ / \pageref{LastPage}}

% Header (expects cpp_logo.png in same folder)
\newcommand{\slideheader}[2]{%
\noindent
\includegraphics[height=0.5in]{cpp_logo.png}%
\hspace{0.3in}%
{\Large \uline{#1 \hfill #2}}%
\par\vspace{0.4in}
}

\begin{document}
\begin{huge}

%%%%%%%%%%%%%%%%%%%%%%%%%%%%%%%%%%%%%%%%%%%%%%%%%%%%%%%%%%%%
% SLIDE 1 - Title Slide
%%%%%%%%%%%%%%%%%%%%%%%%%%%%%%%%%%%%%%%%%%%%%%%%%%%%%%%%%%%%
\slideheader{ECE 5110 Numerical Modeling}{Jacky Li}

\vspace{1.0in}
\begin{center}
{\Huge \textbf{Polynomial Interpolation Seminar}}

\vspace{0.35in}
{\LARGE Unit 2: Interpolation and Polynomial Approximation}\\
\vspace{0.15in}
{\LARGE \textbf{Diode I--V Characteristic Reconstruction}}

\vspace{0.35in}
{\Large February 23, 2025}
\end{center}

\newpage

%%%%%%%%%%%%%%%%%%%%%%%%%%%%%%%%%%%%%%%%%%%%%%%%%%%%%%%%%%%%
% SLIDE 2 - Outline
%%%%%%%%%%%%%%%%%%%%%%%%%%%%%%%%%%%%%%%%%%%%%%%%%%%%%%%%%%%%
\slideheader{Outline}{Interpolation}

\leftskip 0.7in \rightskip 0.7in
\vspace{0.2in}

\begin{itemize}
  \item Introduction: Why Interpolation?
  \item Theory: Lagrange and Newton's Divided Differences
  \item Problem: Diode I--V Curve from Measured Data
  \item Implementation in Python
  \item Results and Error Analysis
  \item Lagrange vs.\ Newton Comparison
  \item Conclusion
\end{itemize}

\leftskip 0in \rightskip 0in
\newpage

%%%%%%%%%%%%%%%%%%%%%%%%%%%%%%%%%%%%%%%%%%%%%%%%%%%%%%%%%%%%
% SLIDE 3 - Why Interpolation?
%%%%%%%%%%%%%%%%%%%%%%%%%%%%%%%%%%%%%%%%%%%%%%%%%%%%%%%%%%%%
\slideheader{Why Interpolation?}{Interpolation}

\leftskip 0.7in \rightskip 0.7in
\vspace{0.15in}

\begin{multicols}{2}
\textbf{The Problem}
\begin{itemize}
  \item In the lab we measure \textbf{discrete data points}
  \item We need to reconstruct the \textbf{continuous curve} between them
  \item Polynomial interpolation: find $P(x)$ of degree $\leq n{-}1$
        passing through all $n$ points
\end{itemize}

\columnbreak

\textbf{Two Classical Methods}
\begin{itemize}
  \item \textbf{Lagrange:} basis-polynomial form
  \item \textbf{Newton:} divided-difference form
  \item Both produce the \textbf{same unique polynomial} (uniqueness theorem)
\end{itemize}
\end{multicols}

\vspace{0.25in}
\textbf{Goal:} Apply both methods to reconstruct a diode I--V curve from measured data and compare.

\leftskip 0in \rightskip 0in
\newpage

%%%%%%%%%%%%%%%%%%%%%%%%%%%%%%%%%%%%%%%%%%%%%%%%%%%%%%%%%%%%
% SLIDE 4 - Theory: Lagrange
%%%%%%%%%%%%%%%%%%%%%%%%%%%%%%%%%%%%%%%%%%%%%%%%%%%%%%%%%%%%
\slideheader{Lagrange Interpolation}{Theory}

\leftskip 0.7in \rightskip 0.7in
\vspace{0.1in}

\textbf{Formula:}
\[
P(x) = \sum_{i=0}^{n-1} y_i \, L_i(x), \qquad
L_i(x) = \prod_{\substack{j=0 \\ j \neq i}}^{n-1}
          \frac{x - x_j}{x_i - x_j}
\]

\vspace{0.15in}
\textbf{Properties:}
\begin{itemize}
  \item Each basis polynomial $L_i(x_j) = \delta_{ij}$ (equals 1 at $x_i$, 0 at all other points)
  \item Explicit closed-form expression
  \item Must recompute \textit{everything} if a new data point is added
\end{itemize}

\leftskip 0in \rightskip 0in
\newpage

%%%%%%%%%%%%%%%%%%%%%%%%%%%%%%%%%%%%%%%%%%%%%%%%%%%%%%%%%%%%
% SLIDE 5 - Theory: Newton
%%%%%%%%%%%%%%%%%%%%%%%%%%%%%%%%%%%%%%%%%%%%%%%%%%%%%%%%%%%%
\slideheader{Newton's Divided-Difference Interpolation}{Theory}

\leftskip 0.7in \rightskip 0.7in
\vspace{0.1in}

\textbf{Formula:}
\[
P(x) = c_0 + c_1(x-x_0) + c_2(x-x_0)(x-x_1) + \cdots
\]

\textbf{Divided differences:}
\[
f[x_i, x_{i+1}] = \frac{f[x_{i+1}] - f[x_i]}{x_{i+1} - x_i},
\qquad c_k = f[x_0, x_1, \ldots, x_k]
\]

\vspace{0.15in}
\textbf{Properties:}
\begin{itemize}
  \item Built incrementally using a divided-difference table
  \item Easy to add new data points (just extend the table)
  \item Generally better numerical stability for large $n$
\end{itemize}

\leftskip 0in \rightskip 0in
\newpage

%%%%%%%%%%%%%%%%%%%%%%%%%%%%%%%%%%%%%%%%%%%%%%%%%%%%%%%%%%%%
% SLIDE 6 - Problem: Diode I-V Curve
%%%%%%%%%%%%%%%%%%%%%%%%%%%%%%%%%%%%%%%%%%%%%%%%%%%%%%%%%%%%
\slideheader{Problem: Diode I--V Curve from Lab Data}{Interpolation}

\leftskip 0.7in \rightskip 0.7in
\vspace{0.05in}

\begin{multicols}{2}

% LEFT
\textbf{Scenario:} Reconstruct a silicon diode's I--V curve from 5 lab
measurements.

\vspace{0.15in}
\textbf{True model} (Shockley):
\[
I_d = I_s \left( e^{\,V_d/(nV_T)} - 1 \right)
\]

We only have \textbf{discrete measured points} with $\pm 2\%$ measurement noise.

\columnbreak

% RIGHT
\begin{center}
\begin{tabular}{cc}
\toprule
$V_d$ [V] & $I_d$ [mA] \\
\midrule
0.30 & $2.31 \times 10^{-6}$ \\
0.50 & $3.97 \times 10^{-4}$ \\
0.65 & 0.0193 \\
0.75 & 0.2590 \\
0.85 & 3.2982 \\
\bottomrule
\end{tabular}

\vspace{0.15in}
{\normalsize Measured diode data ($\pm 2\%$ noise)}
\end{center}
\end{multicols}

\leftskip 0in \rightskip 0in
\newpage

%%%%%%%%%%%%%%%%%%%%%%%%%%%%%%%%%%%%%%%%%%%%%%%%%%%%%%%%%%%%
% SLIDE 7 - Implementation: Lagrange
%%%%%%%%%%%%%%%%%%%%%%%%%%%%%%%%%%%%%%%%%%%%%%%%%%%%%%%%%%%%
\slideheader{Implementation: Lagrange Interpolation}{Jacky Li}

\leftskip 0.7in \rightskip 0.7in
\vspace{0.05in}

\begin{lstlisting}
def lagrange_interpolate(self, x, y):
    n = len(x)
    P = np.zeros(n)
    for i in range(n):
        Li = np.array([1.0])
        denom = 1.0
        for j in range(n):
            if i != j:
                Li = np.convolve(Li, np.array([1.0, -x[j]]))
                denom *= (x[i] - x[j])
        P = P + Li * (y[i] / denom)
    return P, 0
\end{lstlisting}

\vspace{0.15in}
\textbf{Note:} No plots or prints in the method --- plotting is done in \texttt{unit02.py}.

\leftskip 0in \rightskip 0in
\newpage

%%%%%%%%%%%%%%%%%%%%%%%%%%%%%%%%%%%%%%%%%%%%%%%%%%%%%%%%%%%%
% SLIDE 8 - Implementation: Newton
%%%%%%%%%%%%%%%%%%%%%%%%%%%%%%%%%%%%%%%%%%%%%%%%%%%%%%%%%%%%
\slideheader{Implementation: Newton's Divided Differences}{Jacky Li}

\leftskip 0.7in \rightskip 0.7in
\vspace{0.05in}

\begin{lstlisting}
def newton_interpolate(self, x, y):
    n = len(x)
    dd = np.copy(y).astype(float)
    coefs = [dd[0]]
    for j in range(1, n):
        for i in range(n-1, j-1, -1):
            dd[i] = (dd[i]-dd[i-1]) / (x[i]-x[i-j])
        coefs.append(dd[j])
    P = np.array([coefs[0]])
    basis = np.array([1.0])
    for k in range(1, n):
        basis = np.convolve(basis, np.array([1.0, -x[k-1]]))
        P = np.polyadd(P, coefs[k] * basis)
    return P, 0
\end{lstlisting}

\leftskip 0in \rightskip 0in
\newpage

%%%%%%%%%%%%%%%%%%%%%%%%%%%%%%%%%%%%%%%%%%%%%%%%%%%%%%%%%%%%
% SLIDE 9 - Results: I-V Reconstruction
%%%%%%%%%%%%%%%%%%%%%%%%%%%%%%%%%%%%%%%%%%%%%%%%%%%%%%%%%%%%
\slideheader{Results: I--V Curve Reconstruction and Error}{Interpolation}

\leftskip 0.7in \rightskip 0.7in
\vspace{0.05in}

\begin{center}
\includegraphics[width=0.92\textwidth]{unit02_interpolation_results.png}
\end{center}

\vspace{0.1in}
\begin{itemize}
  \item \textbf{Left:} Both Lagrange (blue) and Newton (red) traces overlap --- same polynomial
  \item \textbf{Right:} Absolute error is smallest near data points, grows at edges
\end{itemize}

\leftskip 0in \rightskip 0in
\newpage

%%%%%%%%%%%%%%%%%%%%%%%%%%%%%%%%%%%%%%%%%%%%%%%%%%%%%%%%%%%%
% SLIDE 10 - Results: Data Trend
%%%%%%%%%%%%%%%%%%%%%%%%%%%%%%%%%%%%%%%%%%%%%%%%%%%%%%%%%%%%
\slideheader{Data Trend: True vs Lagrange vs Newton}{Interpolation}

\leftskip 0.7in \rightskip 0.7in
\vspace{0.05in}

\begin{center}
\includegraphics[width=0.62\textwidth]{unit02_data_trend.png}
\end{center}

\vspace{0.15in}
\begin{itemize}
  \item Test voltages \textbf{not in} the original measured data
  \item All three values match closely within the interpolation range
  \item Both methods give \textbf{identical} results (same polynomial confirmed)
\end{itemize}

\leftskip 0in \rightskip 0in
\newpage

%%%%%%%%%%%%%%%%%%%%%%%%%%%%%%%%%%%%%%%%%%%%%%%%%%%%%%%%%%%%
% SLIDE 11 - Comparison Table
%%%%%%%%%%%%%%%%%%%%%%%%%%%%%%%%%%%%%%%%%%%%%%%%%%%%%%%%%%%%
\slideheader{Lagrange vs Newton: Comparison}{Interpolation}

\leftskip 0.7in \rightskip 0.7in
\vspace{0.15in}

\begin{center}
\begin{tabular}{lcc}
\toprule
\textbf{Property} & \textbf{Lagrange} & \textbf{Newton} \\
\midrule
Produces unique polynomial & Yes & Yes \\
Easy to add new data points & No (recompute all) & Yes (extend table) \\
Computational form & Basis polynomials & Divided differences \\
Numerical stability & May suffer for large $n$ & Generally better \\
Coefficient difference & \multicolumn{2}{c}{$\sim 2.22 \times 10^{-16}$ (machine $\varepsilon$)} \\
\bottomrule
\end{tabular}
\end{center}

\vspace{0.2in}
\textbf{Key insight:} Both produce the \textit{same} curve, but Newton is more flexible for incremental data.

\leftskip 0in \rightskip 0in
\newpage

%%%%%%%%%%%%%%%%%%%%%%%%%%%%%%%%%%%%%%%%%%%%%%%%%%%%%%%%%%%%
% SLIDE 12 - Discussion + Conclusion
%%%%%%%%%%%%%%%%%%%%%%%%%%%%%%%%%%%%%%%%%%%%%%%%%%%%%%%%%%%%
\slideheader{Discussion and Conclusion}{Jacky Li}

\leftskip 0.7in \rightskip 0.7in
\vspace{0.05in}

\begin{multicols}{2}
\textbf{Practical Considerations}
\begin{itemize}
  \item Newton's form is more efficient for incremental data collection (e.g., during a lab experiment)
  \item With only $n=5$ points, the 4th-degree polynomial is reasonable \textit{within} the data range
  \item Should \textbf{not} extrapolate beyond data range
  \item For highly nonlinear curves, more points or specialized fitting may be preferred
\end{itemize}

\columnbreak

\textbf{Conclusion}
\begin{itemize}
  \item Reconstructed diode I--V curve from 5 measured points
  \item Both methods yield the \textbf{same} 4th-degree polynomial (uniqueness theorem)
  \item Close agreement with true Shockley model within data range
  \item Newton's divided-difference form is preferred for incremental data
  \item Lagrange provides a more intuitive closed-form expression
\end{itemize}
\end{multicols}

\leftskip 0in \rightskip 0in

\end{huge}
\end{document}
